\section{Introduction}
\label{sec:Introduction}
Scientific workflows execute many tasks whose memory requirements vary
with their inputs and processing behavior. Workflow engines must request
memory before these tasks run. An insufficient reservation can cause an
out-of-memory (OOM) failure and a retry; an excessive reservation occupies
capacity that could serve other tasks. Predictors based on task metadata
and completed executions seek to balance these
costs~\cite{witt2019feedback,lehmann2024ponder,bader2024sizey}. Their practical
value depends on both allocation safety and the resources reserved to
achieve it.

Prelaunch information only partially describes memory demand. A task may
read a small fraction of a large staged index, or access file-backed data
through memory-mapped page faults rather than explicit read calls.
Conversely, a streaming task may process many bytes while retaining little
data in memory. Consumption measurements therefore provide contextual
evidence rather than a direct measure of peak resident memory. The pending
task's consumption is also unavailable at allocation time; only completed
observations can inform its request.

Our prior WORKS paper introduced Consumption-Aware Memory Prediction
(\toolname)~\cite{camp2026works}. It combines static task features and
historical peak memory with consumption measurements from audited tasks.
An auxiliary model predicts consumption before launch; quantile models
and residual calibration then produce a memory request. Task-attributed
eBPF measurements and completed execution outcomes supply feedback for
later allocations. The present paper builds on this architecture to
study a practical deployment question: whether consumption feedback can
reduce wasted memory while maintaining reliable execution as workflows
encounter previously unseen inputs.

The earlier CAMP study compares prediction with and without consumption
features and evaluates selected baselines, ablations, and selective-auditing
policies~\cite{camp2026works}. Here, new Bowtie2 executions distinguish
repeated-input learning from adaptation to a different input study.
A cold execution starts from an offline-trained bundle with no Bowtie2
history; a subsequent warm pass repeats the original inputs. A new-input
stage then continues learning on a disjoint ENA study. To examine the value
of the extra warm pass, a second branch starts the same new inputs directly
from the saved cold checkpoint. These branches share their cold state,
but neither freezes learning during the new workload.

Three considerations guide this evaluation. First, the value of
consumption feedback depends on both measurement and prediction. Explicit
reads and file-backed page-fault activity capture different access behavior;
historical consumption and predicted consumption supply different model
inputs. Separating these effects requires fixed comparison conditions.

Second, auditing incurs execution cost, while its observations can inform only
later tasks. Selecting informative tasks is therefore insufficient by
itself: the resulting allocation savings must be considered alongside the
cost of collecting their measurements. Coverage and reservation cost must
also be assessed together, since fewer underallocations may result from
larger requests. Development-only tuning and a common prelaunch information
boundary keep the component studies comparable.

Third, repeated execution and adaptation to new inputs test different
properties. Repeating a workload supplies additional observations of
familiar inputs; applying the resulting allocator to a different study
also changes task behavior. Comparing branches from a shared checkpoint
can examine the value of the additional pass, while tracking online updates
separately from pre-existing knowledge. Initial-request coverage must be
reported alongside actual grants, failures, and retries: an execution
allowance can prevent an OOM even when the model request is below peak RSS.

Our evaluation addresses three aspects of practical memory allocation:
\begin{itemize}
\item \emph{Measurement and features:} eBPF task filtering and Nextflow
integration, explicit reads alone versus reads with file-backed page-fault
activity, and four feature variants that add historical consumption,
predicted consumption, or both to static features and memory history;
\item \emph{Allocation benefit and cost:} downstream allocations under
selective auditing, estimated memory-time savings relative to auditing
overhead, with coverage and reservation cost evaluated together; and
\item \emph{Online deployment and adaptation:} new Bowtie2 cold-to-warm executions, adaptation to a disjoint input study,
and a shared-checkpoint comparison with and without the extra warm pass.
\end{itemize}
CAMP's architecture and previously reported results remain prior work.
The extension examines when its additional information is useful and
whether the allocation benefit justifies the measurement cost.

Section~\ref{sec:problem} states the allocation problem and information
constraints. Section~\ref{sec:design} summarizes CAMP and its relationship
to WORKS. Sections~\ref{sec:experiments} and~\ref{sec:results} describe the
experimental methods and results, followed by related work
(Section~\ref{sec:related}) and conclusions
(Section~\ref{sec:conclusion}).
